\section{Basic Electromagnetics and Three-Phase Circuit Relations}

\subsection{Magnetic Circuits, Ampere's Law and Flux Linkage ($\oint \mathbf{H}\cdot\mathrm{d}\mathbf{l}$, $N I$, $\Phi$, $\lambda$)}

\subsubsection{Ampere's Circuital Law and Magnetic Circuit MMF Summation ($\sum H_i \ell_i = \sum N I$)}
Ampere's circuital law states that the line integral of magnetic field strength $\mathbf{H}$ around any closed contour equals the total enclosed electric current:
\[
\oint \mathbf{H}\cdot\mathrm{d}\mathbf{l} = \sum I_{\mathrm{enclosed}}.
\]
For a lumped magnetic circuit with discrete segments of length $\ell_i$, cross-sectional area $A_i$, and magnetic field strength $H_i$:
\[
\sum_i H_i\ell_i = \sum F_{\mathrm{source}} = \sum (N I_c)_{\mathrm{source}}.
\]
Each segment presents a magnetic reluctance $\mathcal{R}_i = \frac{\ell_i}{\mu_0 \mu_{r,i} A_i}$, giving Hopkinson's law (magnetic Ohm's law):
\[
\mathcal{F} = \sum (N I_c)_{\mathrm{source}} = \Phi \sum_i \mathcal{R}_i.
\]
In high-permeability iron ($\mu_r \gg 1$), iron reluctance is often negligible compared to the air-gap reluctance $\mathcal{R}_g = \frac{g}{\mu_0 A_g}$.

\subsubsection{Magnetic Flux Density $\mathbf{B}$, Core Area $A_{\mathrm{core}}$, and Total Flux ($\Phi = B_{\mathrm{core}} A_{\mathrm{core}}$)}
Magnetic flux $\Phi$ is the surface integral of flux density $\mathbf{B}$ across cross-section $A_{\mathrm{core}}$:
\[
\Phi = \iint_{A_{\mathrm{core}}} \mathbf{B}\cdot\mathrm{d}\mathbf{A}.
\]
When $\mathbf{B}$ is uniform and perpendicular to the cross-sectional area $A_{\mathrm{core}}$:
\[
\Phi = B_{\mathrm{core}} A_{\mathrm{core}}.
\]

\subsubsection{Faraday's Law of Electromagnetic Induction and Induced EMF ($e = -\mathrm{d}\lambda/\mathrm{d}t$, $E_{\mathrm{RMS}} = 4.44 f N \Phi_m$)}
The total flux linkage $\lambda$ of a coil having $N_{\mathrm{coil}}$ turns enclosing flux $\Phi$ is:
\[
\lambda = N_{\mathrm{coil}}\Phi.
\]
Faraday's law of induction gives the induced electromotive force (EMF):
\[
e = -\frac{\mathrm{d}\lambda}{\mathrm{d}t} = -N_{\mathrm{coil}}\frac{\mathrm{d}\Phi}{\mathrm{d}t}.
\]
The negative sign follows Lenz's law (induced EMF opposes the change in flux).
For a sinusoidal flux $\Phi(t) = \Phi_m \sin(\omega_e t)$:
\[
e(t) = -N_{\mathrm{coil}}\omega_e \Phi_m \cos(\omega_e t)
\implies
E_{\mathrm{RMS}} = \frac{2\pi f N_{\mathrm{coil}}\Phi_m}{\sqrt{2}} \approx 4.44\,f N_{\mathrm{coil}}\Phi_m.
\]

\subsection{Conductor Current Density and Resistance ($J_{\mathrm{cond}}$, $\sigma_{\mathrm{Cu}}$, $R_{\mathrm{cond}}$, Thermal Limits)}

\subsubsection{Conductor Current Density Calculation ($J_{\mathrm{cond}} = I_c / A_{\mathrm{cond}}$)}
The current density $J_{\mathrm{cond}}$ in a single conductor carrying current $I_c$ with cross-sectional area $A_{\mathrm{cond}}$ is:
\[
J_{\mathrm{cond}} = \frac{I_c}{A_{\mathrm{cond}}}\quad[\mathrm{A/mm^2}\text{ or }\mathrm{A/m^2}].
\]
If the phase winding has $a$ parallel circuit paths, the single conductor current is $I_c = I_{\mathrm{ph}} / a$.
Typical machine design limits:
\begin{itemize}
  \item Totally enclosed, naturally cooled: $J_{\mathrm{cond}} \approx 3.0\text{--}4.5\,\mathrm{A/mm^2}$.
  \item Fan-ventilated / air-cooled: $J_{\mathrm{cond}} \approx 4.5\text{--}6.5\,\mathrm{A/mm^2}$.
  \item Direct liquid/hydrogen cooled: $J_{\mathrm{cond}} > 8.0\,\mathrm{A/mm^2}$.
\end{itemize}

\subsubsection{Conductor Resistance and Temperature Dependence ($R_{\mathrm{cond}} = \ell_{\mathrm{cond}} / (\sigma_{\mathrm{Cu}} A_{\mathrm{cond}})$)}
The DC resistance of a single conductor of length $\ell_{\mathrm{cond}}$ is:
\[
R_{\mathrm{cond}} = \frac{\ell_{\mathrm{cond}}}{\sigma_{\mathrm{Cu}} A_{\mathrm{cond}}} = \rho_{\mathrm{Cu}}\frac{\ell_{\mathrm{cond}}}{A_{\mathrm{cond}}},
\]
where $\sigma_{\mathrm{Cu}} \approx 5.8\times 10^7\,\mathrm{S/m}$ at \SI{20}{\degreeCelsius} ($\rho_{\mathrm{Cu}} \approx 1.72\times 10^{-8}\,\Omega\cdot\mathrm{m}$).
At operating temperature $T$:
\[
R_T = R_{20}\left[1 + \alpha_{20}(T - 20^\circ\mathrm{C})\right],
\qquad
\alpha_{20} \approx 0.00393\,\mathrm{K^{-1}}.
\]
At standard hot operating temperature of \SI{75}{\degreeCelsius}, $\sigma_{\mathrm{Cu}} \approx 4.5\times 10^7\,\mathrm{S/m}$.

\subsection{Coil Inductance and Air-Gap Magnetic Field ($L_{\mathrm{coil}}$, $B_g$, $\Phi_g$, $g$)}

\subsubsection{Coil Self-Inductance Definition ($L_{\mathrm{coil}} = \lambda / I_c$)}
The self-inductance $L_{\mathrm{coil}}$ of any coil relates flux linkage to exciting current:
\[
L_{\mathrm{coil}} = \frac{\lambda}{I_c} = \frac{N_{\mathrm{coil}}\Phi}{I_c}\quad[\mathrm{H}].
\]

\subsubsection{Ideal Two-Gap Core Inductance and Air-Gap Flux Density ($B_g = \mu_0 N_{\mathrm{coil}} I_c / (2g)$)}
For a magnetic circuit with infinite core permeability ($\mu_{\mathrm{iron}}\to\infty$), negligible fringing, and two identical series air gaps of physical length $g$ and cross-sectional area $A_g$:
\[
\oint \mathbf{H}\cdot\mathrm{d}\mathbf{l} = 2 g H_g = N_{\mathrm{coil}} I_c
\implies
H_g = \frac{N_{\mathrm{coil}} I_c}{2g}.
\]
The air-gap magnetic flux density $B_g$ is:
\[
B_g = \mu_0 H_g = \frac{\mu_0 N_{\mathrm{coil}} I_c}{2g}.
\]
The total magnetic flux $\Phi_g$ crossing the gap is:
\[
\Phi_g = B_g A_g = \frac{\mu_0 N_{\mathrm{coil}} I_c A_g}{2g}.
\]
The self-inductance of the coil is therefore:
\[
L_{\mathrm{coil}} = \frac{N_{\mathrm{coil}}\Phi_g}{I_c} = \frac{\mu_0 N_{\mathrm{coil}}^2 A_g}{2g}.
\]

\subsection{Exam Recipe: Air-Gap Flux Density $B_g$, Core Flux $\Phi_g$, and Inductance $L_{\mathrm{coil}}$ from Coil MMF}
\textbf{Use when:} Problem provides coil turns $N_{\mathrm{coil}}$, current $I_c$, air-gap length $g$, air-gap pole area $A_g$, and assumes infinite core permeability ($\mu\to\infty$).

\textbf{Steps:}
\begin{enumerate}
  \item Identify the total number of series air gaps crossed by the closed flux line (typically 2 gaps of length $g$, so total gap length is $2g$).
  \item Apply Ampere's law across the gaps: $2g H_g = N_{\mathrm{coil}} I_c \implies H_g = \frac{N_{\mathrm{coil}} I_c}{2g}$.
  \item Compute air-gap flux density: $B_g = \mu_0 H_g = \frac{\mu_0 N_{\mathrm{coil}} I_c}{2g}$.
  \item Calculate pole flux: $\Phi_g = B_g A_g$.
  \item Calculate coil self-inductance: $L_{\mathrm{coil}} = \frac{N_{\mathrm{coil}}\Phi_g}{I_c} = \frac{\mu_0 N_{\mathrm{coil}}^2 A_g}{2g}$.
\end{enumerate}

\textbf{Fast checks:}
\begin{itemize}
  \item Ensure $g$ is converted from $\mathrm{mm}$ to $\mathrm{m}$ ($1\,\mathrm{mm} = 10^{-3}\,\mathrm{m}$).
  \item Ensure $A_g$ is in $\mathrm{m^2}$ ($1\,\mathrm{cm^2} = 10^{-4}\,\mathrm{m^2}$, $1\,\mathrm{mm^2} = 10^{-6}\,\mathrm{m^2}$).
  \item Air-gap flux density $B_g$ in un-saturated electrical machines typically lies in the range $0.5\text{--}1.2\,\mathrm{T}$.
\end{itemize}

\textbf{Common traps:} Forgetting that flux lines must cross the air gap twice to complete a closed loop (factor of 2 in $2g$); mixing area units.

\subsection{Three-Phase AC Power, Star/Delta Configurations and Power Factor ($S_{3\phi}$, $V_{\mathrm{LL}}$, $V_{\mathrm{ph}}$, $I_{\mathrm{line}}$, $I_{\mathrm{ph}}$, $\cos\phi$, $\eta$)}

\subsubsection{Three-Phase Apparent, Active and Reactive Power ($S_{3\phi} = \sqrt{3}V_{\mathrm{LL}}I_{\mathrm{line}}$, $P = S_{3\phi}\cos\phi$)}
Total three-phase apparent power $S_{3\phi}$ in terms of phase and line quantities:
\[
S_{3\phi} = 3 V_{\mathrm{ph}} I_{\mathrm{ph}} = \sqrt{3} V_{\mathrm{LL}} I_{\mathrm{line}}\quad[\mathrm{VA}].
\]
Active power $P$ and reactive power $Q$:
\[
P = S_{3\phi} \cos\phi = 3 V_{\mathrm{ph}} I_{\mathrm{ph}} \cos\phi = \sqrt{3} V_{\mathrm{LL}} I_{\mathrm{line}} \cos\phi\quad[\mathrm{W}],
\]
\[
Q = S_{3\phi} \sin\phi = 3 V_{\mathrm{ph}} I_{\mathrm{ph}} \sin\phi = \sqrt{3} V_{\mathrm{LL}} I_{\mathrm{line}} \sin\phi\quad[\mathrm{var}].
\]

\subsubsection{Motor vs Generator Power Flow and Efficiency ($\eta = P_{\mathrm{mech}} / P_{\mathrm{el}}$)}
\textbf{Motor operation} (electrical power in, mechanical shaft power out):
\[
P_{\mathrm{el}} = \frac{P_{\mathrm{mech}}}{\eta}
\implies
S_{3\phi} = \frac{P_{\mathrm{mech}}}{\eta \cos\phi}.
\]
\textbf{Generator operation} (mechanical shaft power in, electrical power out):
\[
S_{3\phi} = \frac{P_{\mathrm{el}}}{\cos\phi},
\qquad
P_{\mathrm{mech}} = \frac{P_{\mathrm{el}}}{\eta}.
\]

\subsubsection{Star Connection (Y) vs Delta Connection ($\Delta$) Voltage and Current Relations}
\textbf{Star (Y) Connection:}
\[
V_{\mathrm{ph}} = \frac{V_{\mathrm{LL}}}{\sqrt{3}},\qquad
I_{\mathrm{ph}} = I_{\mathrm{line}} = \frac{S_{3\phi}}{\sqrt{3} V_{\mathrm{LL}}}.
\]
\textbf{Delta ($\Delta$) Connection:}
\[
V_{\mathrm{ph}} = V_{\mathrm{LL}},\qquad
I_{\mathrm{ph}} = \frac{I_{\mathrm{line}}}{\sqrt{3}} = \frac{S_{3\phi}}{3 V_{\mathrm{LL}}}.
\]

\subsection{Rotational Speeds, Supply Frequency, Pole-Pairs and Slip ($f$, $p$, $2p$, $n_s$, $n_m$, $\omega_{ms}$, $\omega_m$, $s$, $T_e$)}

\subsubsection{Synchronous Speed Formulas in r/min and r/s ($n_s = 60f/p = f/p$)}
The synchronous speed of the rotating magnetic field depends strictly on electrical frequency $f$ and number of pole-pairs $p$:
\[
n_s = \frac{60 f}{p}\quad[\mathrm{r/min}],
\qquad
n_s = \frac{f}{p}\quad[\mathrm{r/s}].
\]
\begin{center}
\begin{tabular}{@{}cccc@{}}
\toprule
Poles $2p$ & Pole-pairs $p$ & Synchronous speed $n_s$ at \SI{50}{\hertz} & Synchronous speed $n_s$ at \SI{60}{\hertz}\\
\midrule
2 & 1 & \SI{3000}{r/min} & \SI{3600}{r/min}\\
4 & 2 & \SI{1500}{r/min} & \SI{1800}{r/min}\\
6 & 3 & \SI{1000}{r/min} & \SI{1200}{r/min}\\
8 & 4 & \SI{750}{r/min}  & \SI{900}{r/min}\\
\bottomrule
\end{tabular}
\end{center}

\subsubsection{Mechanical Angular Speed ($\omega_m, \omega_{ms}$) vs Electrical Frequency ($\omega_e = 2\pi f$)}
Electrical angular supply frequency:
\[
\omega_e = 2\pi f\quad[\mathrm{rad/s}].
\]
Synchronous mechanical angular speed of the stator air-gap rotating field:
\[
\omega_{ms} = \frac{\omega_e}{p} = \frac{2\pi n_s}{60}\quad[\mathrm{rad/s}].
\]
Actual rotor mechanical rotational angular speed:
\[
\omega_m = \frac{2\pi n_m}{60}\quad(n_m\text{ in r/min}),
\qquad
\omega_m = 2\pi n_m\quad(n_m\text{ in r/s}).
\]

\subsubsection{Induction Machine Slip ($s = (n_s - n_m)/n_s$) and Mechanical Shaft Torque ($T_e = P_{\mathrm{mech}}/\omega_m$)}
Rotor slip $s$ measures the relative velocity between synchronous field and rotor:
\[
s = \frac{n_s - n_m}{n_s} = \frac{\omega_{ms} - \omega_m}{\omega_{ms}}
\implies
\omega_m = \omega_{ms}(1 - s).
\]
Rotor electrical slip frequency:
\[
f_2 = s f,
\qquad
\omega_2 = s \omega_e = p(\omega_{ms} - \omega_m).
\]
Electromagnetic shaft torque $T_e$ developed by the machine:
\[
T_e = \frac{P_{\mathrm{mech}}}{\omega_m} = \frac{P_{\mathrm{gap}}}{\omega_{ms}}\quad[\mathrm{N\cdot m}],
\]
where $P_{\mathrm{gap}}$ is the air-gap power transferred across the gap from stator to rotor.

\subsection{Common Exam Traps and Fast Sanity Checks in Electromagnetics and Power}
\begin{itemize}
  \item \textbf{Poles vs Pole-Pairs:} The formula $n_s = 60f/p$ uses pole-pairs $p$. If problem states ``a 4-pole machine'', then $2p = 4 \implies p = 2$.
  \item \textbf{Speed Units:} Check whether speed is given in $\mathrm{r/min}$ or $\mathrm{r/s}$. Omitting the factor of 60 causes an error of 6000\%.
  \item \textbf{Line vs Phase Voltage in Star:} In star connection, $V_{\mathrm{ph}} = V_{\mathrm{LL}}/\sqrt{3}$. For $400\,\mathrm{V}$ line, $V_{\mathrm{ph}} = 230.9\,\mathrm{V}$.
  \item \textbf{Line vs Phase Current in Delta:} In delta connection, $I_{\mathrm{line}} = \sqrt{3} I_{\mathrm{ph}}$. Do not mix up the $\sqrt{3}$ location!
  \item \textbf{Electrical vs Mechanical Radian Speed:} $\omega_e = p \omega_{ms}$. Always divide $\omega_e$ by $p$ when calculating mechanical angular speed $\omega_{ms}$.
  \item \textbf{Motor Efficiency Location:} For a motor, electrical input is greater than mechanical shaft output ($P_{\mathrm{el}} = P_{\mathrm{mech}}/\eta$). Dividing by $\eta$ increases required VA.
\end{itemize}
